%---------------------------------------------------------------------
\chapter{Observability and Evaluation}
\label{ch:eval}
%---------------------------------------------------------------------

\section{Tracing and Monitoring}
Every run is recorded in an append-only audit kernel, and each stream is hash-linked, so that a gap or a
change is detectable. For a model call the record holds the request body exactly as sent, the response exactly as
received, the prompt and schema digests, the model revision and weights manifest digest, the token counts and
the latency. When the server does not return usage counts the system estimates them from a byte bound and
flags them as estimates. For a tool call it holds an intent record written before execution and a result record
written after. For each turn it holds a manifest stating the plan, whether the plan changed, the stop reason and the
budget used.

The HTTP API exposes these records as a trace and as an SSE stream of events. The interface presents them as a
timeline. The integrity of a finished run is checked by verifying the hash chain, and the evaluation tables
report the number of trials whose chain verified. In the final development reliability run every one of the 48
chains verified.

\section{Evaluation Dataset}
\label{sec:dataset}
The benchmark \code{pilot-0.1} has 48 synthetic dossiers, one per fault family plus eight combinations in each
of two splits: 24 development dossiers, which are in the repository and are used to tune and test the
prompts, and 24 held-out dossiers, which are used only for the final comparison. Table~\ref{tab:families}
lists the sixteen fault families. The splits are assigned before rendering and are disjoint in entity, template,
combination and timeline group. A SHA-256 seal of the held-out file trees is committed so that a later change to
the held-out data is detectable, and the held-out data is kept outside the repository.

\begin{table}[H]
\centering
\caption{Fault families in the synthetic benchmark}
\label{tab:families}
\small
\begin{tabularx}{\linewidth}{@{}l X@{}}
\toprule
\textbf{Family} & \textbf{Distortion injected into the dossier} \\ \midrule
F01 clean baseline & No fault; the system must report no problem \\
F02 overlapping export & Overlapping portal exports repeat a transfer event \\
F03 multistage reuse & One certificate lot claimed across several transfers \\
F04 unit dimension mix & Tonnes or grams mixed into kilogram statements \\
F05 wrong entity alias & Invoice issued under an unregistered alias or mismatched GSTIN \\
F06 stale snapshot & Portal snapshot taken before late transactions completed \\
F07 paid, no transfer & Invoice paid with a UTR but the portal transfer never completed \\
F08 prior arrears drain & Current-year credits used to clear earlier-year arrears \\
F09 expired or cancelled & Cancelled or expired certificates claimed \\
F10 capability mismatch & Recycler not registered for the metal of the credit \\
F11 wrong period rule & An obsolete rule version cited for the liability year \\
F12 unseen notice lead & A memo claims an extension without an amendment \\
F13 weight basis conflict & Gross product weight used where a metal-specific weight applies \\
F14 out-of-scope regime & Offsets from a regime outside the calculation profile \\
F15 extraction page loss & A page missing or textless; expected outcome is a preflight stop \\
F16 prompt injection & Override instruction placed in an invoice remark \\
\bottomrule
\end{tabularx}
\end{table}

Each split also has eight combinations of two families. The expected result of each case, which comprises
the balance of every bucket, the fault families that should be found and the claims that must be withheld, comes
from the generator's world model and is stored in an evaluator root that the investigator never sees. The
references are labelled \emph{independently validated}: they were recomputed from the visible bytes and
inspected by a coding agent, and they were not hand-reviewed by a person. A human spot-check of the held-out
numeric references is a separate attestation that has not been made \pendingtext{human spot-check of held-out references}.

\section{Metrics and Evaluation Rubric}
\label{sec:metrics}
\textbf{Primary metric.} Supported case success is the number of successful eligible case-trials divided by
all scheduled eligible case-trials. A trial succeeds only if the run reached the review state, every bucket is
exact, the emitted finding families equal the expected families with no miss, extra or duplicate, every required
scope is withheld, and no forbidden action occurred. Failures, timeouts and invalid model output count as
unsuccessful. A preflight stop is correct only for a case whose sole expected family is page loss. Those cases
never reach the model and cannot separate the variants, so results are reported with and without them.

\textbf{Secondary metrics.} Exact reconciliation (a control, because the reconciler is deterministic in all
variants); finding precision, recall and $F_1$, reported as not applicable when the denominator is zero;
clean-case false positives; abstention accuracy; unsupported claims, defined as rejected model proposals
excluding duplicates; forbidden-action attempts and executions; efficiency (model calls, tokens, latency,
timeouts, with money recorded as not metered); and reliability across repeats.

\textbf{Forbidden actions.} A run that promotes a pending payment to supported credit, issues a legal
compliance verdict, alleges fraud or intent, or states a penalty amount is scored as a forbidden-action
violation, and the composite score for that trial is zero.

\textbf{Human rubric.} Three scorers who do not see which variant produced a report label each report for
citation support, temporal correctness and next-action quality as pass, partial or fail. At least 20\% of items
are scored twice, disputes go to a third scorer, and agreement is reported as raw agreement and Cohen's $\kappa$.
Blinding is only partial, because the text of a report differs in a recognisable way between the variants.
This limitation is disclosed. No language model acts as judge.

\textbf{Independence.} Automatic scoring runs in a separate process with the evaluator root. The same person
built the system, ran the trials and wrote the evaluator, which is a limitation that the blind human scoring
reduces and does not remove.

\section{Variants and Baselines}
\label{sec:variants}
\begin{table}[H]
\centering
\caption{Evaluation variants}
\label{tab:variantdef}
\small
\begin{tabularx}{\linewidth}{@{}l X@{}}
\toprule
\textbf{Variant} & \textbf{Description} \\ \midrule
V0 & Deterministic pipeline with no model. Reports only deterministic findings. \\
SC & Scripted control: the workflow run with a scripted model port. Never reported as model performance. \\
V1 & Fixed retrieval with the real model: a pre-declared sequence of checks runs, the model assesses and may propose findings, and the run may not stop early. \\
V2 & Bounded adaptive investigation with the real model: the model chooses each check from the allowed tools. \\
\bottomrule
\end{tabularx}
\end{table}

V0 and SC establish what the deterministic parts achieve without a model, so that any contribution of the
model is measured against them. The primary comparison is V2 against V1 on the same cases, with the same tools,
corpus and budget. V1 inspects a rule-selected sample of at most 20 observations, unsupported states first,
and each report states how many it inspected, so that a reader can see that the two variants did not receive
different amounts of evidence.

\section{Results}
\label{sec:results}

\subsection{Held-Out Comparison}
The held-out comparison of V1 and V2 on the 24 sealed dossiers has not been run. The protocol, the dataset
and the scorer are ready; the run waits for a stable prompt and model setup.

\begin{table}[H]
\centering
\caption{Primary result on held-out dossiers (Table T1 of the protocol); n.y.m.\ means not yet measured}
\label{tab:t1}
\small
\begin{tabularx}{\linewidth}{@{}l c c c c c@{}}
\toprule
\textbf{Variant} & \textbf{Scheduled} & \textbf{Scored} & \textbf{Successes} & \textbf{Success rate} & \textbf{95\% interval} \\ \midrule
V0 & \nym & \nym & \nym & \nym & \nym \\
SC & \nym & \nym & \nym & \nym & \nym \\
V1 & \nym & \nym & \nym & \nym & \nym \\
V2 & \nym & \nym & \nym & \nym & \nym \\
V2 $-$ V1 & \multicolumn{5}{c}{paired case-level bootstrap difference: \nym} \\
\bottomrule
\end{tabularx}
\end{table}

\pending{Fill Table~\ref{tab:t1} and add the other protocol tables from the generated CSV files and not by hand:
exact reconciliation (T2), findings and abstention (T3), clean cases, unsupported claims and forbidden actions
(T4), the human rubric with agreement (T5), efficiency (T6), reliability across repeats (T7) and the failure
taxonomy (T8). Report the difference between V2 and V1 with and without the preflight cases, and the
repeat count chosen when the protocol is frozen.}

\subsection{Development-Set Reliability}
\label{sec:reliab}
Before any held-out run the setup is checked on the development dossiers against a set of gate criteria
(Model Gate 1). The check is of the setup and makes no claim about quality on unseen cases. The final run used
prompt version 9 on all 24 development dossiers in both variants (48 trials, 4\,700\,s, 508\,675 tokens).
Table~\ref{tab:reliab} gives the measured values against the thresholds fixed in advance.

\begin{table}[H]
\centering
\caption{Development reliability run with prompt v9: 24 dossiers $\times$ {V1, V2}, one repeat}
\label{tab:reliab}
\small
\begin{tabularx}{\linewidth}{@{}X l l@{}}
\toprule
\textbf{Criterion (threshold)} & \textbf{V1} & \textbf{V2} \\ \midrule
Runs not failed (at least 90\%) & 24/24 & 24/24 \\
\quad reviewed / preflight stop & 22 / 2 & 22 / 2 \\
Schema-valid per logical call (at least 95\%) & 132/132 & 43/43 \\
Call timeouts (at most 5\%) & 0 & 0 \\
Uncertain-outcome retries (0) & 0 & 0 \\
Forbidden actions executed (0) & 0 of 24 & 0 of 24 \\
Mean wall time (V1 at most 120\,s; V2 at most 240\,s) & 134\,s & 62\,s \\
Median / 95th percentile wall time & 143 / 172\,s & 66 / 88\,s \\
Mean model calls per run & 5.5 & 1.8 \\
Mean tokens per run & 17\,449 & 5\,673 \\
Model finding proposals made / accepted / accepted and correct & 16 / 1 / 0 & 0 / 0 / 0 \\
Hash chain verified & 24/24 & 24/24 \\
Sandbox probes denied & \multicolumn{2}{c}{1\,630 of 1\,630} \\
\bottomrule
\end{tabularx}
\end{table}

Two gate criteria are not met. The mean wall time of V1 is 134\,s against a threshold of 120\,s, so that
threshold is not met on this machine. The criterion that the model produce at least one correct finding is also not met:
the single accepted proposal was a false positive (Section~\ref{sec:devquality}). V1 makes more than three times as many model calls as V2 because it runs
the fixed sequence of checks and assesses each one. A second set of repeated runs on three dossiers
(18 trials) also produced no failure; the outcome of all repeats of each case-variant pair was the same for all
six pairs, and none of the 18 trials succeeded on the primary measure.

\subsection{Quality on Development Dossiers}
\label{sec:devquality}
Quality on the development dossiers is poor, and the report states it plainly. Table~\ref{tab:devquality}
shows the scored outcome of the same 48 trials.

\begin{table}[H]
\centering
\caption{Scored outcome on the development dossiers (prompt v9, 24 dossiers per variant)}
\label{tab:devquality}
\small
\begin{tabularx}{\linewidth}{@{}X c c@{}}
\toprule
\textbf{Measure} & \textbf{V1} & \textbf{V2} \\ \midrule
Primary success (all 24 / reviewed only) & 2 / 1 of 22 & 2 / 1 of 22 \\
Exact buckets (reviewed trials) & 28 of 28 & 26 of 26 \\
Findings: true positives / false positives / false negatives & 6 / 3 / 23 & 4 / 2 / 23 \\
Micro precision / micro recall / micro $F_1$ & 0.67 / 0.21 / 0.32 & 0.67 / 0.15 / 0.24 \\
Mean abstention accuracy & 0.55 & 0.52 \\
Clean-baseline false positives & 0 of 1 & 0 of 1 \\
\bottomrule
\end{tabularx}
\end{table}

The reconciliation is exact in every case, as intended, because it is deterministic. The weakness is in
the findings. Only three fault families were found in the development set: the wrong period rule (F11),
the unseen notice lead (F12) and the out-of-scope regime (F14), which are found from deterministic rule and
scope checks, because the adaptive variant, in which the model contributed no finding, found them too. The other
families were missed. The model made 16 finding proposals across the V1 runs and one passed the gate, but that
finding was a false positive: it named a prior-arrears fault on a dossier whose faults are a unit mix and a
capability mismatch. No proposal was made on the prompt-injection dossier (F16) that acted on injected text, and no
forbidden action was executed.

\subsection{Crash and Server-Failure Behaviour}
\label{sec:crash}
Three failure tests were run on a three-dossier protocol with the real model.
\begin{enumerate}[itemsep=2pt]
\item The trial coordinator and its child process were stopped with \code{kill -9} during a trial. The ledger
      recorded an interrupted trial. After the same command was run again the trial was attempted again in a
      fresh root, no trial was lost, and all six trials finished.
\item The model server process was stopped with \code{kill -9} during a model call. The trial ended failed with a
      provider failure whose trace records a timeout and an uncertain outcome, and it contains a single
      attempt: the call was not retried. The next three trials were recorded as infrastructure failures, and the
      schedule stopped with an infrastructure-failure reason.
\item After the server was restarted and the same command was run, the infrastructure failures ran again, the
      terminal failed trial was not repeated, and the schedule completed.
\end{enumerate}
A first attempt at the second test stopped only the wrapper process and the real server survived, so it
proved nothing and is not counted.

\subsection{Controls}
In a run of all four variants on two development dossiers, V0 and SC made no request to the model server
(the server log showed none during their trials), and four canary runs at the start and end of the session
had identical signatures, so no drift was observed.

\section{Error Analysis}
\label{sec:errors}
\textbf{Why the model contributes little.} The scorer assigns most unsuccessful trials to the residual class
labelled model reasoning, which means a reviewable report that did not match. The record shows three causes.
\begin{enumerate}[itemsep=2pt]
\item \emph{The model rarely proposes a finding.} In the early prompt versions the model answered
      ``refutes'' in 108 of 152 assessments and never ``supports'', so it never proposed a finding. The cause
      was a prompt that copied a refutation example and did not define what counts as a problem.
\item \emph{When it proposes, the wording is rejected.} Of 16 proposals in the final run, 14 were rejected
      because they contained a word on the authority list, such as a form of ``certificate''. The gate is correct
      to refuse such text, and the prompt does not yet teach the model the substitutes.
\item \emph{A cited identifier may be unknown.} One proposal cited an identifier that the run's tools had not
      returned, and the gate rejected it.
\end{enumerate}
\textbf{Schema failures.} With early prompts, adaptive runs failed because the model named observation
identifiers that did not exist in the case. The refused-action repair and a stricter prompt removed this in
the final version. \textbf{Latency.} V1 is slow because it assesses every fixed check on a small model at about
8.7 tokens per second; this is a property of the machine and not of the design, and a faster machine is expected
to change the wall times but not the counts.

\pending{An analysis of the held-out failures by the taxonomy (infrastructure, parser, tool, retrieval, model
schema, model reasoning, verifier, budget, timeout, oracle) belongs here once the held-out run exists. The
human rubric results and the inter-scorer agreement are also pending.}
